\documentclass[../../../../uem_mat_mei_q.tex]{subfiles}

\begin{document}
    $n$を2以上の自然数とする．整数$k$に対して，%
    座標平面上に点$\mathrm{P}_k$を次のようにとる．

    \myspaceb%
    $\displaystyle%
        \mathrm{P}_k\left(\cos\myfracb{2k\pi}{n},\;\sin\myfracb{2k\pi}{n}\right)%
    $\\
    点$\mathrm{P}_{k-1}$と点$\mathrm{P}_k$を結ぶ線分の長さを$L_k$とする．このとき，

    \myspaceb%
    $\displaystyle%
        L_k=\myboxb{キ}%
    $\\
    である．これより，

    \myspaceb%
    $\displaystyle%
        \lim_{n \to \infty}\sum_{k=1}^nL_k=\myboxb{ク}%
    $\\
    がわかる．

    \vspace{.5\baselineskip}
    \textsf{キ，クの解答群}
    
    {\renewcommand{\arraystretch}{1.7}%
        \vspace{-.2\baselineskip}%
        \hspace{\zw}%
        \begin{tabularx}{\dimexpr\linewidth-2\zw}{@{}LLL@{}}
            ⓪\; $\myfracb{2\pi}{n}$ & ①\; $2\sin\myfracb{\pi}{n}$ &%
            ②\; $2\cos\myfracb{2\pi}{n}$ \\
            ③\; $\tan\myfraca{\pi}{2n}$ & ④\; $2\tan\myfraca{\pi}{2n}$ & ⑤\; $\pi$ \\
            ⑥\; $\myfraca{3\pi}{2}$ & ⑦\; $2\pi$ & ⑧\; $4\pi$ \\
            ⑨\; $\infty$
        \end{tabularx}%
    }
\end{document}
